\documentclass[10pt]{article}

\usepackage[utf8]{inputenc}
\usepackage[T1]{fontenc}
\usepackage{amsmath,amssymb}
\usepackage{booktabs}
\usepackage{enumitem}
\usepackage[margin=0.78in]{geometry}
\usepackage{hyperref}

\hypersetup{colorlinks=true,linkcolor=blue,urlcolor=blue}
\setlist{itemsep=2pt,topsep=3pt,parsep=0pt,partopsep=0pt}
\setlength{\parindent}{0pt}
\setlength{\parskip}{3pt}

\newcommand{\R}{\mathbb{R}}
\newcommand{\dd}{\,\mathrm{d}}
\newcommand{\norm}[1]{\left\lVert#1\right\rVert}

\title{Summer Bootcamp -- Session 2\\Hands-On Calculus and ODE Exercises}
\author{Nathan Sauldubois}
\date{August 2026}

\begin{document}
\maketitle

\section{Python and Elementary Financial Calculations}

This first part requires only short Python functions and elementary algebra.

\paragraph{Exercise 1 -- Coupon-bond yield, duration, and repricing.}
A bond pays $4$ at $t=1$, $4$ at $t=2$, and $104$ at $t=3$. Its market
price is $96.50$. We denote by $B(y)$ the price of the bond as a function of
the continuously compounded yield $y$. Thus
\[
B(y)=4e^{-y}+4e^{-2y}+104e^{-3y}.
\]
\begin{enumerate}
\item write Python functions for $B(y)$, $B'(y)$, and $B''(y)$;
\item find the yield with bisection on an interval whose endpoint prices
bracket $96.50$;
\item verify the result with Newton's method;
\item compute modified duration and convexity,
\[
D_{\mathrm{mod}}=-\frac{B'(y)}{B(y)},
\qquad
\mathcal C=\frac{B''(y)}{B(y)};
\]
\item estimate the new price after yield moves of $\pm25$ and $\pm100$ basis
points using duration alone and then duration--convexity;
\item compare both approximations with exact repricing.
\end{enumerate}

\section{Explicit Derivatives}

\paragraph{Exercise 2 -- One-dimensional derivatives.}
Differentiate and simplify each function. State its domain.
\[
f_1(x)=\sin(3x^2-1),
\qquad
f_2(x)=x^2e^{-3x},
\]
\[
f_3(x)=\frac{\log x}{1+x^2},
\qquad
f_4(x)=\log(1+e^x),
\qquad
f_5(x)=\sqrt{1+x^2}.
\]
For $f_2$ and $f_4$, compute the second derivative as well.

\paragraph{Exercise 3 -- Gradient, Hessian, and Jacobian.}
Let
\[
f(x,y)=2x^2-xy+3y^2+4x,
\qquad
g(x,y)=\begin{pmatrix}e^x\cos y\\x-y^2\end{pmatrix}.
\]
Using the method of your choice, compute symbolically:
\begin{enumerate}
\item the gradient $\nabla f(x,y)$;
\item the Hessian $H_f(x,y)$;
\item the Jacobian $J_g(x,y)$.
\end{enumerate}

\paragraph{Exercise 4 -- Multivariate chain rule by two methods.}
Define
\[
q(u,v)=ue^v,
\qquad
\Phi(x,y)=(x+y^2,xe^y),
\qquad F=q\circ\Phi.
\]
\begin{enumerate}
\item write $F(x,y)$ explicitly and compute $\partial_xF$ and $\partial_yF$;
\item compute $J_\Phi(x,y)$ and $\nabla q(u,v)$;
\item recover $\nabla F$ from
$\nabla F=J_\Phi^\top\nabla q(\Phi)$;
\item write the increment calculation
$DF(x,y)h=Dq(\Phi(x,y))D\Phi(x,y)h$ and verify that it gives the same
coefficients of $h_1$ and $h_2$.
\end{enumerate}

\section{Explicit Integration}

\paragraph{Exercise 5 -- One-dimensional integrals.}
Evaluate analytically and state the method used:
\[
I_1=\int_0^1(3x^2-2x+1)\dd x,
\qquad
I_2=\int_0^T(1+t)^{-r}\dd t,
\]
\[
I_3=\int_0^1xe^{x^2}\dd x,
\qquad
I_4=\int_0^1x\log x\dd x,
\qquad
I_5=\int_0^1\frac{1}{1+x^2}\dd x.
\]
For $I_2$, distinguish $r=1$ from $r\neq1$ and recover the exceptional case
by a limit. For $I_5$, use the substitution $x=\tan u$.

\paragraph{Exercise 6 -- The Gaussian integral.}
Let
\[
G=\int_{-\infty}^{\infty}e^{-x^2}\dd x.
\]
The objective is to obtain the value of $G$ by a formal calculation.
\begin{enumerate}
\item write $G^2$ as an integral over $\R^2$ using independent variables
$x$ and $y$;
\item change to polar coordinates, stating the domain and the Jacobian;
\item evaluate the radial and angular integrals and deduce $G$;
\item use the result to compute
$\int_{-\infty}^{\infty}e^{-x^2/(2\sigma^2)}\dd x$ for $\sigma>0$;
\item explain how this calculation determines the normalizing constant of a
Gaussian density.
\end{enumerate}

\paragraph{Exercise 7 -- Multiple integrals and change of variables.}
Compute
\[
J_1=\int_{-1}^2\int_0^3(2x-y)\dd y\dd x,
\]
\[
J_2=\iint_D(x+2y)\dd x\dd y,
\qquad
D=\{(x,y):x\geq0,\ y\geq0,\ x+2y\leq2\},
\]
and
\[
J_3(R)=\iint_{1\leq x^2+y^2\leq R^2}(x^2+y^2)^2\dd x\dd y,
\qquad R>1.
\]
\begin{enumerate}
\item compute $J_1$ in both integration orders;
\item write both possible iterated-integral descriptions of $D$ and compute
$J_2$;
\item transform $J_3(R)$ to polar coordinates, including the Jacobian $r$;
\item specialize the result to $R=3$;
\item check the three answers with a short numerical quadrature in Python.
\end{enumerate}

\section{Black--Scholes Greeks}

For a non-dividend-paying asset,
\[
C=S\Phi(d_1)-Ke^{-rT}\Phi(d_2),
\qquad
d_1=\frac{\log(S/K)+(r+\tfrac12\sigma^2)T}{\sigma\sqrt T},
\qquad d_2=d_1-\sigma\sqrt T.
\]
Here $\Phi$ is the cumulative distribution function of a standard normal
random variable:
\[
\Phi(x)=\frac{1}{\sqrt{2\pi}}
\int_{-\infty}^{x}e^{-u^2/2}\dd u,
\qquad
\phi(x)=\Phi'(x)=\frac{1}{\sqrt{2\pi}}e^{-x^2/2}.
\]

\paragraph{Exercise 8 -- Derive the Greeks.}
The closed-form formulas for the Greeks are deliberately not supplied: they
must be obtained by differentiating the Black--Scholes price.

With all other variables held fixed, the Greeks are defined by the following
partial derivatives:
\[
\boxed{
\Delta=\frac{\partial C}{\partial S},\qquad
\Gamma=\frac{\partial^2 C}{\partial S^2},\qquad
\mathrm{Vega}=\frac{\partial C}{\partial\sigma}.}
\]
\[
\boxed{
\Theta=-\frac{\partial C}{\partial T},\qquad
\mathrm{Rho}=\frac{\partial C}{\partial r}.}
\]
Here $T$ denotes time to maturity, so calendar-time Theta is
$-\partial C/\partial T$.

\begin{enumerate}
\setcounter{enumi}{1}
\item derive a useful identity relating $S\phi(d_1)$ and
$Ke^{-rT}\phi(d_2)$;
\item use the chain rule to derive the closed-form formulas for all five
Greeks, simplifying with the identity from the previous question;
\item Derive the short-maturity at-the-money approximation as follows.
\begin{enumerate}[label=(\alph*)]
\item Use the at-the-money-forward convention
\[
K=Se^{rT},
\]
and explain why it is equivalent to equality between the spot price $S$ and
the discounted strike $Ke^{-rT}$.
\item Substitute this relation into $d_1$ and $d_2$ and prove that
\[
d_1=\frac{\sigma\sqrt T}{2},
\qquad
d_2=-\frac{\sigma\sqrt T}{2}.
\]
\item Rewrite the call price in the form
\[
C=S\left[
\Phi\left(\frac{\sigma\sqrt T}{2}\right)
-\Phi\left(-\frac{\sigma\sqrt T}{2}\right)
\right],
\]
then use the symmetry $\Phi(-z)=1-\Phi(z)$.
\item Using the expansion
\[
\Phi(z)=\frac12+\frac{z}{\sqrt{2\pi}}+O(z^3)
\qquad (z\to0),
\]
show that, as $T\downarrow0$,
\[
C=S\sigma\sqrt{\frac{T}{2\pi}}+O(T^{3/2}).
\]
\item Interpret the leading term: describe separately the effects of $S$,
$\sigma$, and $T$, and explain why the short-maturity price is of order
$\sqrt T$ rather than order $T$.
\end{enumerate}
\end{enumerate}

\section{Numerical Schemes for ODEs}

For $y'=f(t,y)$ on $t_n=nh$, recall
\[
y_{n+1}=y_n+hf(t_n,y_n)
\quad\text{and}\quad
y_{n+1}=y_n+hf(t_{n+1},y_{n+1}).
\]

\paragraph{Exercise 9 -- Euler accuracy and the classical inequalities.}
For $y'=-2y$, $y(0)=1$, on $[0,3]$:
\begin{enumerate}
\item derive explicit and implicit Euler updates and their closed forms;
\item implement both schemes for $h=0.5,0.25,0.125,0.0625$;
\item compute
$\max_n|y(t_n)-y_n|$ using $y(t)=e^{-2t}$ and estimate the convergence order;
\item derive the explicit stability inequality
$|1-2h|<1$ and the corresponding restriction on $h$;
\item verify that the implicit update multiplier $1/(1+2h)$ has modulus
below one for every $h>0$;
\item explain why stability does not imply accuracy for a very large step.
\end{enumerate}

\section{Solving Scalar and Multidimensional ODEs}

\paragraph{Exercise 10 -- Scalar linear equations.}
Solve exactly and verify by differentiation:
\[
y'+4y=2t,\qquad y(0)=-1,
\]
\[
y'-\frac{1}{1+t}y=(1+t)^2,\qquad y(0)=2,
\]
\[
y'-y=e^{2t},\qquad y(0)=0.
\]
Use variation of the constant for the first and third equations and an
integrating factor for the second.

\paragraph{Exercise 11 -- Nonlinear scalar equations.}
For an autonomous ODE $y'=f(y)$, an \emph{equilibrium} is a constant solution
$y(t)\equiv y^\star$, equivalently a value satisfying $f(y^\star)=0$.

Solve the following equations analytically whenever possible:
\[
y'=y(1-y),\qquad y(0)=y_0\in(0,1),
\]
\[
y'=-y^3,\qquad y(0)=1,
\]
\[
y'=t(1+y^2),\qquad y(0)=0.
\]
\begin{enumerate}
\item For $y'=y(1-y)$, first identify the equilibrium solutions. For a
non-equilibrium solution, divide by $y(1-y)$ and use
\[
\frac{1}{y(1-y)}=\frac1y+\frac1{1-y}.
\]
Integrate both sides, determine the integration constant from $y(0)=y_0$,
then solve explicitly for $y(t)$.
\item For $y'=-y^3$, rewrite the equation with all powers of $y$ on the
left. Integrate, use $y(0)=1$, and choose the branch consistent with the
initial condition. Verify the answer by differentiation.
\item For $y'=t(1+y^2)$, separate the variables and recognize an elementary
antiderivative on the left. After applying $y(0)=0$, solve for $y(t)$ and
inspect where that expression ceases to be finite.
\item For each solution, determine its maximal interval of existence. Do not
infer this interval only from the ODE: inspect the explicit denominator or
the inverse function introduced during the calculation.
\item Draw the phase line for the two autonomous equations and determine
the stability of every equilibrium.
\end{enumerate}

\paragraph{Exercise 12 -- Linear systems in dimension two.}
Consider
\[
x'(t)=Ax(t),
\qquad
A=\begin{pmatrix}-2&1\\0&-4\end{pmatrix},
\qquad
x(0)=\begin{pmatrix}1\\2\end{pmatrix}.
\]
\begin{enumerate}
\item compute the eigenvalues and eigenvectors of $A$;
\item decompose $x(0)$ into eigenvectors and obtain the exact solution;
\item verify the solution using $x(t)=e^{tA}x(0)$;
\item implement explicit and implicit Euler in vector form;
\item compare both numerical trajectories with the exact solution.
\end{enumerate}

\paragraph{Exercise 13 -- Forced multidimensional system.}
Let
\[
x'=Ax+b,
\qquad
A=\begin{pmatrix}-2&0\\1&-1\end{pmatrix},
\qquad
b=\begin{pmatrix}2\\0\end{pmatrix},
\qquad x(0)=\begin{pmatrix}0\\2\end{pmatrix}.
\]
\begin{enumerate}
\item compute the equilibrium $x^\star=-A^{-1}b$;
\item use
$x(t)=x^\star+e^{tA}(x(0)-x^\star)$ to obtain an exact solution;
\item derive the same answer component by component;
\item implement Euler schemes and verify convergence to $x^\star$.
\end{enumerate}

\paragraph{Exercise 14 -- Second-order equation as a system.}
Solve
\[
y''+4y'+3y=0,
\qquad y(0)=0,
\qquad y'(0)=2.
\]
\begin{enumerate}
\item use the characteristic polynomial;
\item rewrite the equation as a first-order system in $(y,y')$;
\item compute the eigenvalues of the companion matrix;
\item implement explicit Euler for the system and compare with the exact
solution obtained in the first question.
\end{enumerate}

\end{document}
